\subsection{ISOMORPHISMS OF WELL-ORDERED SETS}

THEOREM 3. Let E and F be two well-ordered sets. Then at least one of the following two statements is true:

\begin{enumerate}
\def\labelenumi{(\arabic{enumi})}
\item
  there exists a unique isomorphism of \(\mathbf{E}\) onto a segment of \(\mathbf{F}\) ;
\item
  there exists a unique isomorphism of \(\mathbf{F}\) onto a segment of \(\mathbf{E}\) .
\end{enumerate}

Let \(\mathfrak{F}\) be the set of mappings of subsets of E into F such that each mapping is defined on a segment of E and is an isomorphism of this segment onto a segment of F. Then the set \(\mathfrak{F}\) , ordered by the relation `` \(v\) extends \(u\) '' between \(u\) and \(v\) , is inductive. For if \(\mathfrak{G}\) is a totally ordered subset of \(\mathfrak{F}\) , the union S of the domains of the mappings \(u \in \mathfrak{G}\) is a union of segments of E and is therefore itself a segment of E. If \(v\) is the least upper bound of \(\mathfrak{G}\) in \(\Phi(E, F)\) (no. 4, Example 2), then \(v(S)\) is the union of the ranges of the mappings \(u \in \mathfrak{G}\) and is therefore a segment of F. Finally, for each pair of elements \(x\) , \(y\) of S such that \(x < y\) there exists \(u \in \mathfrak{G}\) whose domain contains both \(x\) and \(y\) (because \(\mathfrak{G}\) is totally ordered); and since \(v(x) = u(x) < u(y) = v(y)\) , \(v\) is an isomorphism of S onto \(v(S)\) , and our assertion is proved.

Now let \(u_0\) be a maximal element of \(\mathfrak{F}\) (no. 4, Theorem 2) and let \(S_0\) be the segment of \(E\) which is the domain of \(u_0\) . If we show that either \(S_0 = E\) or \(u_0(S_0) = F\) , the theorem will be proved. Let us argue by contradiction and suppose that \(S_0 \neq E\) and \(u_0(S_0) \neq F\) . There will then be an element \(a \in E\) and an element \(b \in F\) such that \(S_0 = ]\leftarrow, a[\) and \(u_0(S_0) = ]\leftarrow, b[\) (no. 1, Proposition 1). Extend \(u_0\) to a mapping \(u_1\) of the segment \(]\leftarrow, a]\) into \(F\) by putting \(u_1(a) = b\) ; since \(u_1\) is an isomorphism of \(]\leftarrow, a]\) onto the segment \(]\leftarrow, b]\) , this contradicts the maximality of \(u_0\) in \(\mathfrak{F}\) .

The uniqueness asserted in Theorem 3 is a consequence of the following Lemma :

LEMMA 4. Let E, F be two well-ordered sets and let f, g be two increasing mappings of E into F such that \(f(\mathbf{E})\) is a segment of F and g is strictly increasing; then \(f(x) \leqslant g(x)\) for all \(x \in E\) .

Suppose, on the contrary, that the set of elements \(y \in E\) such that \(f(y) > g(y)\) is not empty; then this set will have a least element \(a\) . If \(x < a\) , we have then \(f(x) \leqslant g(x) < g(a) < f(a)\) since \(g\) is strictly increasing. Since \(f(E)\) is a segment of \(F\) , there exists \(z \in E\) such that \(g(a) = f(z)\) ; \(f\) is increasing, so that \(f(z) < f(a)\) implies \(z < a\) . Hence

\[
f (z) \leqslant g (z) <   g (a) = f (z),
\]

which is absurd.

COROLLARY 1. The only isomorphism of a well-ordered set E onto a segment of E is the identity mapping of E onto itself.

Put \(\mathbf{F} = \mathbf{E}\) in Theorem 3.

COROLLARY 2. Let E, F be two well-ordered sets. If there exists an isomorphism f of E onto a segment T of F and an isomorphism g of F onto a segment S of E, then we must have S = E, T = F, and g, f are inverses of each other.

For \(g \circ f\) is an isomorphism of \(\mathbf{E}\) onto the segment \(g(\mathbf{T}) \subset \mathbf{S}\) of \(\mathbf{E}\) ; by Corollary 1 we have \(g(\mathbf{T}) = \mathbf{S} = \mathbf{E}\) , and \(g \circ f\) is the identity mapping of \(\mathbf{E}\) . Similarly, \(f \circ g\) is the identity mapping of \(\mathbf{F}\) , whence the result.

COROLLARY 3. Every subset A of a well-ordered set E is isomorphic to a segment of E.

By virtue of Theorem 3 it is enough to prove that there exists no isomorphism g of E onto a segment of A of the form \(S_{a}\) . If there were, g would then be a strictly increasing mapping of E into E such that \(g(a) \in S_{a}\) , in other words such that \(g(a) < a\) ; but this inequality contradicts Lemma 4 (with f as the identity mapping).

